\begin{table}[htbp]
\centering
\begin{tabular}{lrrrrrr}
\toprule
 & \multicolumn{2}{c}{Chain Ladder} & \multicolumn{2}{c}{Hierarchical} & \multicolumn{2}{c}{Cape Cod} \\
\cmidrule(lr){2-3} \cmidrule(lr){4-5} \cmidrule(lr){6-7}
Origin & Reserve & LR & Reserve & LR & Reserve & LR \\
\midrule
Origin 1 & 0 & 0.7205 & 0 & 0.7205 & 0 & 0.7205 \\
Origin 2 & 15,126 & 0.7165 & 15,159 & 0.7165 & 15,209 & 0.7165 \\
Origin 3 & 26,257 & 0.7375 & 26,000 & 0.7375 & 25,619 & 0.7375 \\
Origin 4 & 34,538 & 0.6943 & 35,075 & 0.6944 & 35,874 & 0.6944 \\
Origin 5 & 85,302 & 0.6797 & 87,294 & 0.6798 & 90,234 & 0.6800 \\
Origin 6 & 156,494 & 0.6717 & 160,597 & 0.6719 & 166,584 & 0.6723 \\
Origin 7 & 286,121 & 0.6451 & 297,878 & 0.6459 & 314,665 & 0.6470 \\
Origin 8 & 449,167 & 0.5983 & 482,247 & 0.6006 & 528,056 & 0.6037 \\
Origin 9 & 1,043,242 & 0.6010 & 1,112,597 & 0.6058 & 1,200,821 & 0.6119 \\
Origin 10 & 3,950,815 & 0.6329 & 4,105,566 & 0.6431 & 4,240,563 & 0.6519 \\
Total & 6,047,061 & 0.6699 & 6,322,415 & 0.6717 & 6,617,625 & 0.6737 \\
\bottomrule
\end{tabular}
\caption{Posterior mean reserves and ultimate loss ratios by origin, using the exact analytic approach.}
\label{tab:saluz_reserve_means_analytic}
\end{table}